\chapter{Economic and environmental implications of onshoring US lithium production}\label{chap:lith}

\begin{center} 
    Ongoing research. 

    Anthony L. Cheng\footnote{Principal researcher and author}\textsuperscript{,}\footnote{Department of Engineering and Public Policy, Carnegie Mellon University \label{ch4-note1}}, Erica R.H. Fuchs\textsuperscript{\ref{ch4-note1}}, Valerie J. Karplus\textsuperscript{\ref{ch4-note1},}\footnote{Scott Institute for Energy Innovation, Carnegie Mellon University}, Jeremy J. Michalek\textsuperscript{\ref{ch4-note1},}\footnote{Department of Mechanical Engineering, Carnegie Mellon University}
\end{center}

\clearpage

\subsection{Abstract}

Rapid growth in batteries for electric vehicles and energy storage has intensified global demand for critical minerals, particularly lithium. Although geological lithium resources are widely distributed, incumbent supply chains are highly dependent on China for refined lithium materials production, raising concerns about supply chain vulnerability and environmental impacts. While most Chinese lithium chemical production relies on imported ores, domestic production from lepidolite feedstocks currently appears to operate on or near the margin of the global supply curve, as inferred from recent price-sensitive entry and exit behavior. Yet, its cost and environmental performance are poorly characterized in the open literature, limiting assessment of whether alternative production pathways can compete economically and environmentally. We develop an integrated, empirically grounded and fully transparent bottom-up technoeconomic and life cycle greenhouse gas emission model for three lithium projects representing distinct resource-pathway combinations: U.S. evaporative brine at Silver Peak, the only active domestic producer; U.S. sulfuric acid-leach clay at Thacker Pass, currently under construction; and Chinese sulfate roasting lepidolite at Jianxiawo, a large near-marginal producer on the global supply curve. We find that resource-pathway pairings create structural cost and emissions differences that persist across geographies and dominate over other location-specific factors such as reagent costs, labor costs, or electricity grid impact factors. Silver Peak (\$5,000 [\$4,100---5,900]/t \ce{Li2CO3}; 5.2 [4.5---5.8] t~\ce{CO2eq}/t) and Thacker Pass (\$8,750 [\$7,400---12,300]/t; 17.0 [14.9---19.0] t \ce{CO2eq}/t) compare favorably to Jianxiawo (\$10,400 [\$9,500---12,300]/t; 22.6 [17.0---33.0] t \ce{CO2eq}/t), where ranges reflect plausible variation in the dominant cost and emissions drivers. These estimates are made under conservative project-specific designs for Thacker Pass and an optimistic ore grade for Lepidolite, suggesting the true gap is likely wider. Contrary to conventional wisdom that Chinese operations benefit from lax environmental regulation, we find that high environmental compliance costs of lepidolite production in China, driven by hazardous waste disposal requirements for roasting residues, further disadvantage the pathway relative to US sources. These results suggest that growth of U.S. lithium production has potential to economically displace marginal lepidolite production in China while offering lower climate change externalities.

\clearpage
\subsection{Introduction}

Lithium is an essential component of nearly all current electric vehicle batteries and large-scale battery energy systems, yet it is projected to experience an increase in demand by a factor of 20 to 40 between 2020 and 2040~\cite{internationalenergyagencyGlobalCriticalMinerals2024, xuFutureMaterialDemand2020, greimAssessmentLithiumCriticality2020}. While materials from battery recycling could dampen this demand, production from raw materials extraction will still need to produce more than 80\% of future demands under most scenarios~\cite{xuFutureMaterialDemand2020}. Natural endowments of lithium are relatively widespread, but current production is geographically concentrated, while also mismatched with demand sources (\ref{fig:04-li_distribution}). Australia, Chile, and Argentina dominate primary extraction while China dominates refining and materials production, and yet most demand for lithium products, such as for electric vehicles, is in China, the U.S., and Europe \cite{chengElectricVehicleBattery2024, greitemeierChinasHoldLithiumion2025}. This concentration has raised concerns about supply chain resilience for countries that rely on lithium for key energy technologies \cite{internationalenergyagencyGlobalCriticalMinerals2024, olivettiLithiumIonBatterySupply2017, chengElectricVehicleBattery2024}. These concerns span both geopolitical risks such as export restrictions and resource nationalism \cite{bremmerRiseFallResource2009, prinaceraiGeographyControlDeep2024} and physical risks from natural disasters, labor strikes, or environmental impacts \cite{helbigSupplyRisksAssociated2018, schrijversReviewMethodsData2020}, leading to an increasing interest in ``reshoring'' and ``friend-shoring'' supply chains. 

\begin{figure}[!h]
    \centering
    \altfig{A series of horizontal stacked bars in one chart showing the geographic distribution of lithium across five supply chain stages: resources, reserves, mining and extraction, refining, EV cathode production, and end demand. Resources and reserves are concentrated in Chile, Australia, and Argentina. Australia and Chile dominate extraction. China dominates refining with a large majority share. End demand for cathodes and EVs is split among China, Europe, the United States, South Korea, and Japan, highlighting the mismatch between where lithium is found, where it is processed, and where it is consumed.}
        \includegraphics[width=\linewidth]{figures/Chapter 4/04-li-distribution.png}
    \altfigend
    \caption{Geographic distribution of lithium (Li) across resources, production, and demand (2023).}
    \small{There is a significant mismatch between the location of lithium resources, reserves, lithium production, lithium refining, and end lithium demand. The following codes are used to represent countries or regions: USA, United States of America; AUS, Australia; CHL, Chile; CAN, Canada; CHN, China; EUR, Europe; ARG, Argentina; BRA, Brazil; ZIM, Zimbabwe; JPN, Japan; KOR, South Korea; Other any other country not explicitly listed here. Data for the year 2023, with resources, reserves, extraction, and the quantity of refined lithium from the USGS Mineral Commodity Survey 2025 \cite{usgeologicalsurveyMineralCommoditySummaries2025}, refining percentages from the IEA Global Critical Minerals Outlook 2024\cite{internationalenergyagencyGlobalCriticalMinerals2024}, and cathode and EV percentages from the IEA Global EV Outlook 2024\cite{internationalenergyagencyGlobalEVOutlook2024}. Note that lithium mining totals in the United States are not disclosed by USGS, but are reported in other sources\cite{howellNevadaHasLoads2024}.}
    \label{fig:04-li_distribution}
\end{figure}

Global lithium production currently draws on two primary resource types: brines and hard rocks, with sedimentary clay deposits representing a significant but underdeveloped third source. Spodumene hard rock deposits, mined primarily in Australia and refined in China, currently represent the largest share of overall lithium production, and sit toward the lower-to-middle of the cost curve at roughly \$5,000---8,000/t LCE~\cite{flemingLongrunAvailabilityMineral2024, wesselkaemperOptimizationProcessingCosts2025}. Continental brine deposits found in desert salars in Chile and Argentina represent the lowest-cost production globally, with high brine grade and high aridity enabling solar evaporation-based concentration at costs as low as \$3,000---5,000/t LCE~\cite{flemingLongrunAvailabilityMineral2024}. However, with the high water usage of these resources in water-stressed regions \cite{sunRobustAssessmentsLithium2025, halkesLifeCycleAssessment2024}, alternate `direct lithium extraction' (DLE) technologies are emerging as a lower-footprint alternative for processing these brine-based resources, though these are currently only operating in China and at relatively limited scale. A range of other `unconventional' resources, such as geothermal brines, jadarite, zinnwaldite, and others, currently contribute negligibly to global supply. 

Of these other resources, sedimentary clay deposits, most prominently found in the western United States, represent the largest undeveloped unconventional lithium source. Largely because of these clays, the United States held about 12\% of world resources in 2023 \ref{fig:04-li_distribution}, which has since increased to 20\% as of 2026 \cite{usgeologicalsurveyMineralCommoditySummaries2026}. However, regulatory timelines, depressed prices, uncertainty in financing, and community and Indigenous opposition have led to delays in developing these resources; currently, only two projects have been permitted: Thacker Pass and Rhyolite Ridge, both clay-based resources. With the latter suffering from a collapse in funding~\cite{haasRhyoliteRidgeLithium2025}, Thacker Pass remains the only U.S. project under construction, with a target operational date of 2028 - two decades after significant exploration activity started in 2007 \cite{alonzoEvenLithiumPrices2025, NolanTestimony2025, calaway2025AnnualReport2025}. However, its target annual production capacity of 160,000 tonnes LCE by 2044~\cite{sgscanadaincSK1300Technical2024} would be roughly 10\% of world production in 2025~\cite{usgeologicalsurveyMineralCommoditySummaries2026}. Given only that the Silver Peak evaporative brine facility in Nevada has been the only lithium producer in the United States since the beginning of the 21st century, with a current nameplate capacity of roughly 6,000 tonnes, this marks a significant increase in US production. In sum, if all Nevada clay projects operated by publicly listed companies were to come online, that could result in a total of 320,000 tonnes LCE of annual production (including Thacker Pass)~\cite{sgscanadaincSK1300Technical2024, fluorenterprisesRhyoliteRidgeLithiumBoron2022, dyerAmendedResourceEstimate, laneNI43101Technical2024, riordanTonopahLithiumClaims2023, samariMineralResourceEstimate2024}. Thus, while other Direct Lithium Extraction sites in the United States are under construction, their relatively limited total annual production capacity of 40,000 tonnes~\cite{worleyltdParadoxLithiumProject, gallardoSouthWestArkansas2025, brushNi43101Technical2023} and relatively low reported operating costs (\$5,600-7,400 per tonne LCE) point towards an initial focus on better understanding U.S. clay resources. 

On the other hand, China's domestic lithium production has grown roughly tenfold between 2017 and 2025, from about 37,000 to an estimated 330,000 tonnes LCE, driven primarily by rapid expansion of low-grade \textit{lepidolite} hard rock mica deposits~\cite{benchmarkteamRestartCATLsJianxiawo2026, usgeologicalsurveyMineralCommoditySummaries2026}. At costs of roughly \$10,000---16,000/t~LCE~\cite{liuChinaLithiumBoom2024, atkinsonLithiumPriceReversion2023, jamasmieCATLCutLithium2024}, or roughly the 90th percentile of the global cost curve~\cite{sulcaLithiumPricesWhere2024}, lepidolite is the source most sensitive to price fluctuations and the relevant marginal producer against which alternative suppliers must compete.

A key question for countries seeking to reduce their dependence on incumbent supply chains is therefore whether domestic production can undercut this marginal cost and compete on the global market. Lepidolite production has contracted during periods of low prices, with disruptions to lepidolite supply from Jiangxi province repeatedly triggering significant volatility in Chinese lithium futures markets \cite{liChinasCATLSuspends2025, benchmarkteamRestartCATLsJianxiawo2026}. While the dynamics for these projects are examined in more detail in Section~\ref{sec:04-jianxiawo}, this price-driven entry and exit behavior, combined with the scale of existing installed capacity, suggests that lepidolite is likely to remain the relevant marginal producer as demand grows, making it an appropriate competitive benchmark.

However, as demand for lithium continues to grow, displacing lepidolite production on the cost curve would require the production of new alternative sources, at a competitive cost. Whether that investment is justified from a financial perspective depends on the expected return relative to risk and alternative uses of capital, typically evaluated via net present value (NPV) or internal rate of return (IRR) metrics \cite{pindyckIrreversibilityUncertaintyInvestment1990, brennanEvaluatingNaturalResource1985}. However, from a competitive market-based supply security perspective, the more relevant threshold is whether a project's average lifetime operating cost is competitive with the current marginal producer: the price at which an incumbent's production becomes uneconomic sets the ceiling that a new entrant must undercut to capture market share \cite{tiltonMineralEconomicsPolicy2016}. It is this long-run cost competitiveness, rather than project-level returns or the exact cost curve of non-marginal production, that determines whether domestic supply development could durably compete on the global scale. 

Production characteristics such as costs and emissions are not separable from the resource itself: the chemistry and geology of the deposit --- whether brine, hard rock, or clay sources --- determine the extraction and processing pathway, which in turn drives both the cost structure and also externalities such as emissions and environmental damages. This means that geographically shifting production is not simply a matter of replicating existing processes, but also that these production-resource pairings could inherently provide advantages or disadvantages. This resource-pathway relationship creates opportunities for policy to shape supply chain configurations in ways that could simultaneously address security, cost, and environmental objectives, rather than trading one off against another \cite{riofrancosSecuritySustainabilityNexus2023, haddadHowMakeLithium2023, buschEffectsDemandRecycling2025}. Some studies \cite{ujiPursuingDecarbonizationNational2023} have identified that communicating the national security and domestic sourcing benefits of decarbonization projects can help navigate local public opposition. The extent to which such `win-win-win' outcomes are achievable depends on the relative cost and emissions profiles of competing pathways across production locations, which this work aims to characterize.

Assessing whether these outcomes are achievable requires characterization of the cost and emissions profiles of these competing pathways. Yet, the literature remains incomplete, particularly for the Chinese lepidolite operations that define the current benchmark for marginal production. While a substantial body of literature has examined models of lithium production through life cycle assessment with varying levels of detail (see the review of over 40 such studies in Appendix \ref{app:04-methods-litreview}), process-based technoeconomic analyses are comparatively sparse (e.g.  \cite{wesselkaemperOptimizationProcessingCosts2025, warrenTechnoEconomicAnalysisLithium2021}), and studies that jointly evaluate both dimensions across multiple resource types are rarer still, with only \citet{mousavinezhadNextLithiumBoom2025, gaoLithiumExtractionHard2023} to our knowledge. \citet{mousavinezhadNextLithiumBoom2025} is the only study to evaluate multiple production pathways along the full mine-to-product scope in both dimensions, but focus solely on four projects representing US-based pathways (U.S. evaporative brine, U.S. direct lithium extraction (DLE) from brine, U.S. clay, and spodumene mined in Canada and processed in the US). On the other hand, lepidolite has received limited attention and highly varied results: \citet{liuCarbonFootprintBatteryGrade2025} estimated its global warming potential (GWP) to be the highest amongst production pathways in China at roughly 24.5 tonnes \ce{CO2eq}/t \ce{Li2CO3}. In contrast, \citet{wangLifeCycleEnvironmental2026} estimated a much higher GWP impact of lepidolite at 66.0 tonnes \ce{CO2eq}/t \ce{Li2CO3}. Its production costs are even more poorly characterized in the open literature, with only \citet{gaoLithiumExtractionHard2023} reporting that the production cost of lithium carbonate from lepidolite concentrate via the sulfate method was about \$4,500 per ton \ce{Li2CO3} - excluding the direct production costs of lepidolite concentrate production from ore, and not providing an open, examinable model. In contrast, news reports and grey literature have identified lepidolite production costs to vary widely between USD\$10,000 and USD\$16,000 \cite{atkinsonLithiumPriceReversion2023, liuChinaLithiumBoom2024, jamasmieCATLCutLithium2024}. We plot these values, in addition to literature estimates of costs and emissions for the Silver Peak and Thacker Pass operations as dashed lines in \ref{fig:04-costemissions}, noting that each of these estimates come from different studies with different scopes and assumptions. 

We aim to assess whether the cost and emissions advantages suggested by resource type hold, and what this implies for the competitiveness of lithium production in the United States relative to lepidolite in China as the current marginal producer.

\subsection{Scope and Methodology}\label{sec:04-scopemethod}

As a first step towards addressing these literature gaps, this work evaluates and compares the production costs and cradle-to-gate greenhouse gas emissions of three lithium projects representing distinct resource-pathway combinations, using a transparent, empirically-grounded technoeconomic-life cycle assessment framework and consistent assumptions across US and Chinese sites. This process-level framework directly interlinks production costs and emissions, providing a static treatment of lithium production but empirical grounding of production, laying the groundwork for more richly represented dynamic analyses, such as responses to price signals, technological learning, policy feedbacks, or geopolitical shifts. Furthermore, given that the emissions literature is substantially richer and more methodologically mature than the technoeconomic literature for lithium, consistency between modeled emissions profiles and reported life cycle assessment values can thus provide an independent cross-check on the underlying process model, and therefore the cost estimates derived from it. 

In order of increasing complexity of mineral resource and processing pathway, we define flowsheets for each project as shown in (\ref{fig:04-step_breakdown}), based on publicly available project disclosures, patents, and other literature. In particular, Silver Peak and Thacker Pass are owned and operated by Albemarle and Lithium Americas, respectively, two companies that are publicly traded on US stock exchanges. As a result, under SEC Regulation S-K 1300\cite{SECgovModernizationProperty}, they are required to disclose mineral resources and reserves through reports written by an independent Qualified Person, from which high-level flowsheets can be determined \cite{srkconsultingSECTechnicalReport2023a, sgscanadaincSK1300Technical2024}. On the other hand, publicly available detailed representations of lepidolite projects in China are limited, given the lack of similar disclosure requirements in China. We focus on Jianxiawo (Yichun, China, \ref{fig:04-step_breakdown}c), the largest lepidolite project in China by production volume: its integration with CATL, the world's largest battery manufacturer, means that its operational changes, including production suspensions and restarts, attract substantial coverage in industry press and commodity market reporting, providing more contextual data than is available for most Chinese lepidolite operations (e.g.~\cite{liChinasCATLSuspends2025,jamasmieCATLCutLithium2024,benchmarkteamRestartCATLsJianxiawo2026,benchmarkteamCATLSuspendsOperations}). Furthermore, its usage of the sulfate roasting pathway, the dominant processing route for lepidolite in China~\cite{dengSystematicReviewSulfate2026}, makes it representative of the pathway, beyond being the largest individual project. These three projects represent key competing pathways relevant to supply chain security; the full set of generic lithium carbonate production pathways is described in Appendix \ref{fig:c-pathways}. 

\begin{figure}
    \centering
    \altfig{Three process flow diagrams showing the step-by-step production sequences for each lithium project modeled. Panel a shows Silver Peak, an evaporative brine operation. Panel b shows Thacker Pass, a clay-based operation. Panel c shows Jianxiawo, a lepidolite hard-rock operation. Each diagram shows production steps as rounded boxes grouped within shaded facility boundaries, with inward arrows indicating key reagent inputs and outward arrows indicating coproducts or waste streams at each step.}
        \includegraphics[width=0.76\linewidth]{figures/Chapter 4/04-stepwise.png}
    \altfigend
    \caption{Process overview diagrams for each project modeled: (a) Silver Peak; (b) Thacker Pass; (c) Jianxiawo}
    \footnotesize{Each subfigure shows the sequence of production steps, key reagents consumed (arrows inward) and coproducts or waste streams produced (arrows outward). Each shaded box in the background represents a facility, where production steps (rounded boxes on top of shaded boxes) are co-located, with the same sitewide operating parameters.}
    \label{fig:04-step_breakdown}
\end{figure} 

The primary methodology we employ in this analysis is a series of process-based technoeconomic models, from which environmental impacts can be estimated and policy interventions can be applied. The modeling framework is structured hierarchically across four levels of resolution, enabling consistent comparison across projects while preserving sensitivity to pathway-specific technology requirements. At the \textit{supply chain} level, facilities are linked by \textit{transportation routes} into a directed network from raw material extraction to refined \ce{Li2CO3} product. Each \textit{facility} aggregates co-located production steps operating under shared geographic and operational parameters (e.g. electricity grid composition, labor costs, utility prices, and plant schedule), ensuring that location-specific cost and emissions factors are applied consistently. Within each facility, \textit{production steps} represent individual unit operations (e.g., acid leaching, impurity removal, lithium carbonate precipitation), each defined by a primary input and output stream, secondary reagent inputs, and coproduct or waste streams routed to defined sinks. At each step, three material volume quantities are tracked: the primary input, the process basis (the internal scaling quantity, such as \si{m^3} slurry processed or tonnes solids crushed), and the primary output, which are related using two conversion factors (from input to step basis and step basis to output) that accommodate the diversity of natural reporting units across feedstock types and process steps. Reagent consumption is calculated as a function of constituent concentration in the process stream, allowing stoichiometric or empirical dosing ratios to be applied directly to the relevant constituent rather than to bulk throughput. Each production step is parameterized using \textit{machine blocks}, which specify equipment sizing, capital cost scaling relationships, utility intensities, and labor requirements. This separation between process chemistry (defined at the production step level) and equipment characteristics (defined at the machine block level) allows alternative equipment configurations or technology variants to be substituted without altering the underlying mass balance, enabling systematic sensitivity and scenario analysis. Costs and emissions are computed at the production step level and aggregated upward; because material flows are propagated from a terminal production target back through the full process chain while taking into account process unit conversions and yield rates, throughput, reagent consumption, and equipment sizing at every step are directly linked to the final output volume. For further detail on this hierarchical modeling system and a formal mathematical formulation, see Appendix \ref{app:04-methods-system}. 

Six overarching facility types are defined, as shown in \ref{fig:04-step_breakdown}: brine evaporation, blast mining, pregnant leach solution (PLS) production from batch chemical treatment, PLS production from sulfuric acid leaching, PLS production from sulfate roasting and water leaching, and lithium carbonate production (which is parameterized differently depending on if the PLS is a chloride or sulfate-based solution). Each of these facilities consist of multiple production steps, for a total of thirty production steps modeled (\ref{fig:04-step_breakdown}). Production volumes are modeled at approximately 6,000~t~\ce{Li2CO3}/yr for Silver Peak~\cite{srkconsultingSECTechnicalReport2023a}, 66,000~t/yr for Thacker Pass (consistent with the prior feasibility disclosure~\cite{rothFeasibilityStudyNational2022}), and 46,000~t~LCE/yr for Jianxiawo (the most commonly cited production figure~\cite{flemingLongrunAvailabilityMineral2024, andrezejewskiResourcesEnergyQuarterly2025}).

For each pathway, we compute average lifetime operating costs of lithium carbonate production in USD per metric tonne \ce{Li2CO3}. As discussed in the introduction, this metric is more appropriate for assessing whether a project can sustainably compete with current or future marginal producers than a full NPV or IRR analysis. For each production step, operating expenditures (OPEX) are computed as the sum of variable costs (materials, labor, utilities) and fixed operating costs (maintenance and overhead), expressed per unit of annual \ce{Li2CO3} production. Capital expenditures (CAPEX), such as equipment, buildings, and auxiliary systems, are estimated using installation cost factors and equipment scaling relationships. For continuous processes, overall equipment costs are scaled with throughput using power law exponents (e.g. the six-tenths rule) while batch processes are expanded by adding parallel lines \cite{petersPlantDesignEconomics1991, ausimmMonograph27Cost, ulrichGuideChemicalEngineering1984}. These capital costs are not included in the primary cost comparison given the focus on operating competitiveness, but are reported for reference and cross-comparison with the literature. 

These results are extended to capture the environmental impacts of production, where impact factors for reagent use, energy consumption, and utilities are multiplied by the mass and energy flows through the system (see Appendix \ref{app:04-methods-system-emissions}). Scope 1 emissions result from specified process emissions (which are limited in our three cases) or direct combustion of fuel, primarily for heating process throughputs. Scope 2 emissions result from electricity usage, which is then multiplied by the local grid emissions factor. Scope 3 emissions are calculated by multiplying each reagent's annual consumption by its cradle-to-gate emissions factor, reflecting the upstream production burden of that reagent as a purchased commodity input. This methodology is applied uniformly across all pathways, with geography-specific emissions factors applied where production location affects this upstream burden (e.g., soda ash from trona mining versus the Solvay process). 

For both costs and emissions, conservative, midpoint, and optimistic values are determined for each machine block design and material input cost, yielding sensitivity ranges that represent plausible operating variation across the dominant uncertainty drivers for each project. Process chemistry (e.g. reagent stoichiometry, conversion factors, and yield rates) is held fixed within each pathway, since these are generally determined by the underlying mineralogy and cannot be varied independently of the process design. Similarly, utility prices, labor rates, and land costs are not varied in the sensitivity analysis, as these are set by location rather than by operational choices, and their cross-pathway differences are instead captured through the location-specific parameters applied at the facility level.

While Silver Peak and Jianxiawo are modeled using flowsheets derived from project-specific disclosures and patents respectively, the clay pathway is represented using a generic acid-leach flowsheet grounded in the fundamental process chemistry of Nevada sedimentary clay deposits rather than the specific design of Thacker Pass. Several other clay projects in Nevada share similar mineralogy and processing fundamentals (e.g. Rhyolite Ridge~\cite{fluorenterprisesRhyoliteRidgeLithiumBoron2022}, Bonnie Claire~\cite{samariMineralResourceEstimate2024}, Tonopah Flats~\cite{dyerAmendedResourceEstimate}, TLC~\cite{riordanFalchaniLithiumProject2024}) and a generic flowsheet provides a more transferable baseline for this class of resource than a project-specific model would. The key distinction is in impurity management: because sulfuric acid dissolves both lithium and magnesium simultaneously from magnesium-rich hectorite, the generic flowsheet precipitates all magnesium with quicklime, whereas the Thacker Pass project-specific design employs a magnesium crystallization route that recycles magnesium compounds as reagents, reducing lime consumption to approximately 2,630~kg/t at the cost of higher thermal energy demand~\cite{sgscanadaincSK1300Technical2024}. This omission renders the generic flowsheet conservative relative to the Thacker Pass design specifically, but representative of what a first-generation clay operation without project-specific optimization would require. The cost and emissions implications of this choice are discussed in Section~\ref{sec:04-thackerpass}.

\subsection{Results and Discussion}
\begin{figure}[!h]
    \centering
    \altfig{Scatter plot with production cost in dollars per tonne of lithium carbonate on the x-axis and life cycle greenhouse gas emissions in tonnes CO2-equivalent per tonne lithium carbonate on the y-axis. Three projects are plotted: Silver Peak at approximately 4,980 dollars per tonne and 5.2 tonnes CO2-eq, Thacker Pass at approximately 8,750 dollars per tonne and 17.0 tonnes CO2-eq, and Jianxiawo at approximately 10,400 dollars per tonne and 22.6 tonnes CO2-eq. Each project has a darker shaded rectangle showing the sensitivity range from this study and a lighter rectangle spanning the range of values reported in the literature, with individual literature data points marked as lines. A dark grey spur extends from the Jianxiawo midpoint showing modeled cost reduction (approximately 900 dollars less) under a hazardous waste non-compliance scenario.}
        \includegraphics[width=\textwidth]{figures/Chapter 4/04-costemissions.png}
    \altfigend
    \caption[Production cost (\$/t \ce{Li2CO3}) versus life cycle greenhouse gas emissions (t \ce{CO2eq}/t \ce{Li2CO3}) for three key lithium projects representing distinct resource-pathway combinations.]{Production cost (\$/t \ce{Li2CO3}) versus life cycle greenhouse gas emissions (t \ce{CO2eq}/t \ce{Li2CO3}) for three key lithium projects representing distinct resource-pathway combinations: Silver Peak (evaporative brine, United States), Thacker Pass (clay, United States), and Jianxiawo (lepidolite, China)}
    \small{Darker shaded rectangle span the range from this study, with the central point indicating the study midpoint; these ranges reflect variation in the three largest sources of cost and emissions uncertainty for each project, identified through sensitivity analysis (\ref{fig:04-silverpeak-sensitivity}, \ref{fig:04-thackerpass-sensitivity}, \ref{fig:04-jianxiawo-sensitivity}). The complete uncertainty range incorporating all parameters simultaneously is shown in \ref{fig:c-extendedresults}; the ranges shown here are intended to represent plausible operating variation rather than absolute worst-case bounds. The dark grey spur extending from the Jianxiawo midpoint shows the modeled cost reduction under a hazardous waste non-compliance scenario, in which lepidolite roasting residues are disposed of as non-hazardous waste rather than under current Chinese regulations (see Section~\ref{sec:04-jianxiawo}). Lighter rectangles span the range of values reported across grey and peer-reviewed literatures, with each line representing one data point. The lighter shaded rectangles span the full range of values reported across grey and academic literatures, with each line representing a data point, with A: \citet{srkconsultingSECTechnicalReport2023a}; B: \citet{sgscanadaincSK1300Technical2024}; C: \citet{ambroseUnderstandingFutureLithium2020a}; D: \citet{mousavinezhadNextLithiumBoom2025}; E: \citet{flemingLongrunAvailabilityMineral2024}; F: \citet{wesselkaemperOptimizationProcessingCosts2025}; G: \citet{chenLithiumSpotPrice2023}; H: \citet{atkinsonLithiumPriceReversion2023}; I: \citet{liuChinaLithiumBoom2024}; J: \citet{liLithiumWeaknessPushes2024}; K: \citet{sunRobustAssessmentsLithium2025}; L: \citet{manjongLifeCycleModelling2021}; M: \citet{ambroseUnderstandingFutureLithium2020b}; N: \citet{schenkerCurrentFutureImpacts2025}; O: \citet{khakmardanLifeCycleAssessment2024}; P: \citet{iyerLifecycleAnalysisLithium2024}; Q: \citet{royLithiumClayAssessing2024}; R: \citet{liuCarbonFootprintBatteryGrade2025}; S: \citet{wangLifeCycleEnvironmental2026}. See \ref{tab:c1-litreview-matrix} for a summary of all literature studies.}
    \label{fig:04-costemissions}
\end{figure}

Our results indicate that US-based lithium production from evaporative brine and sulfuric acid-based clay pathways have lower production costs and life cycle emissions than production from Chinese-based lepidolite sulfate roasting pathways. The Silver Peak project occupies the lowest-cost, lowest-emission frontier across both our estimates and the literature, at an estimated cost of about \$5,000 [\$4,100-\$5,900] per tonne LCE, and 5.2 [4.5-5.8] tonnes \ce{CO2eq} per tonne LCE. Modeled production costs and emissions from the clays in the generic version of the Thacker Pass flowsheet are closer to modeled production from lepidolite from Jianxiawo, with lower midpoint costs (\$8,750 [\$7,400-\$12,300] vs. \$10,400 [\$9,500-\$12,300]) and emissions (17.0 [14.9-19.0] tonnes \ce{CO2eq} vs. 22.6 [17.0-33.0] tonnes \ce{CO2eq}). 

Our modeled ranges for all three projects are broadly consistent with reported literature values (\ref{fig:04-costemissions}), confirming the cross-check between emissions and cost model structure described in Section~\ref{sec:04-scopemethod}. The substantially wider ranges for lepidolite across both the literature and our study also reflects the data scarcity in reported production costs and emissions for this pathway, underscoring the importance of better understanding this major source of lithium production.

From this, we see that Silver Peak's cost and emissions advantage is robust across all modeled ranges, but the ranges aiming to represent plausible operating variation for Thacker Pass and Jianxiawo overlap, particularly when taking into account a scenario where hazardous wastes generated by Jianxiawo's production process are treated as if they were non-hazardous. This indicates that the relative competitiveness of U.S. clay projects versus Chinese lepidolite is directionally favorable, but this is sensitive to project-specific assumptions explored in the sections below. 

\subsubsection{Results for the Silver Peak Project}\label{sec:04-silverpeak}

The Silver Peak project, located in a historic silver mining town in southwestern Nevada, is the only active operating lithium producing facility within the United States. Currently owned by Albemarle, it has produced lithium since 1966 through a traditional evaporative brine, batched-chemical process, with a current production target of roughly 6000 tonnes of lithium carbonate each year \cite{srkconsultingSECTechnicalReport2023a}. In this pathway, lithium-bearing brine is pumped from subsurface aquifers and concentrated through solar evaporation in large, shallow ponds before downstream chemical purification yields battery-grade \ce{Li2CO3} (\ref{fig:04-step_breakdown}a). Most desert salar brines are composed of sodium chloride solutions with lithium, magnesium, calcium and potassium ions \cite{garrettHandbookLithiumNatural2004}; as water evaporates from brines, many of these ions naturally precipitate from the solution. However, magnesium and lithium are typically soluble, and magnesium removal is typically the most chemically demanding purification step for lithium brines\cite{torresSalarAtacamaLithium2024, flemingLongrunAvailabilityMineral2024}. As opposed to South American salars such as Salar de Atacama ($\sim$2,000 mg/L), Silver Peak's brine resource is relatively dilute at approximately \qty{95}{\mg\per\L}, requiring extensive evaporation before chemical processing can begin; in our model, approximately 95\% of pumped water evaporates as the brine concentrates to 6,750 PPM strong brine. However, Silver Peak's relatively low magnesium content allows for an all-chemicals, batched precipitation process with lime and soda ash, a simpler production method than solvent extraction or ion exchange processes necessary for other salar brines~\cite{srkconsultingSECTechnicalReport2023a, flexerLithiumRecoveryBrines2018}. Thus, in spite of this resource disadvantage, Silver Peak's cost and emissions profile is the most favorable of the three projects modeled here, a consequence of its structural reliance on solar evaporation for lithium concentration and said low magnesium concentration, shifting its cost and emissions profile towards maintenance and upkeep of the solar evaporation system rather than the more chemically intensive ore-based production pathways.

Evaluating the step-by-step costs breakdown (\ref{fig:04-silverpeak-summary}a), we see that three steps dominate: evaporation ponds ($\sim$\$1,350/t), batch treatment ($\sim$\$1,300/t), and brine pumping ($\sim$\$1,000/t). The evaporation pond cost is predominantly labor and the "Other" category, here represented primarily by maintenance operations on the large pond footprint. Pond-level pumping electricity is assumed to be included with the pumping step, another major cost contributor. On the other hand, lithium carbonate precipitation is the primary step with significant material costs, making soda ash the primary reagent cost driver. The overall cost structure is labor- and capital-intensive rather than reagent-intensive, reflecting an operating cost that is relatively insensitive to global markets and more sensitive to local operating conditions.

From an emissions perspective, brine pumping ($\sim$2.0 t \ce{CO2eq}/t), batch chemical treatment ($\sim$1.9 t \ce{CO2eq}/t), and lithium carbonate precipitation ($\sim$1.1 t \ce{CO2eq}/t) are by far the dominant emissions steps (\ref{fig:04-silverpeak-summary}b). The former is Scope 2, driven by electricity consumption for pumping using an emissions factor based on the Nevada grid\cite{usepaSummaryData2025}, while the latter two are Scope 3, from embodied emissions in lime and soda ash production. While soda ash production in the United States is generally cleaner ($\sim$0.3-0.7 t \ce{CO2eq}/t \ce{Na2CO3}) due to its direct production from Trona rather than the heat-intensive Solvay process used in China ($\sim$0.7-1.2 t \ce{CO2eq}/t \ce{Na2CO3}) \cite{wangGreenhouseGasesRegulated2024, scottCanSyntheticSoda2023, royLithiumClayAssessing2024}, the heavy use of lime for magnesium removal ($\sim$0.83-1.7 t \ce{CO2eq}/t \ce{CaO} at 1.5 tonnes \ce{CaO} per tonne \ce{Li2CO3}) and soda ash to remove the calcium introduced by the lime and precipitate lithium carbonate (2.4 tonnes \ce{Na2CO3} per tonne \ce{Li2CO3}) results in significant impacts from the embodied emissions associated with the production of these upstream inputs. However, this emissions profile does also indicate that pumping is by far the most significant electricity consumer, and that decarbonization of the Nevada grid would have a direct and substantial impact on Silver Peak's emissions intensity.

\begin{figure}
    \centering
    \altfig{Two-panel stacked bar chart for the Silver Peak brine project. Panel a shows production cost in dollars per tonne lithium carbonate broken down by process step. Panel b shows greenhouse gas emissions in tonnes CO2-equivalent per tonne lithium carbonate broken down by the same steps. Error bars on each step show the sensitivity envelope from varying machine design parameters and material input costs between conservative and optimistic values.}
        \includegraphics[width=1\linewidth]{figures/Chapter 4/04-Silver Peak-Steps-Summary.png}
    \altfigend
    \caption{Costs(a) and emissions(b) for each step of the Silver Peak Project.}
    \small{Error bars represent a sensitivity envelope formed by simultaneously applying conservative and optimistic values for all machine design parameters and material input costs that carry defined scenario variation; they do not reflect variation in process chemistry, reagent stoichiometry, or steps for which no scenario variation is defined (e.g., tailings management). As a result, the envelope is wider for steps with more varied parameters and narrower or absent for steps with few or none. Error bars should therefore be interpreted as indicative of the dominant sources of operating uncertainty rather than as statistical confidence intervals.}
    \label{fig:04-silverpeak-summary}
\end{figure}

\begin{figure}
    \centering
    \altfig{Two-panel tornado chart for the Silver Peak project. Panel a shows the top five inputs with the greatest influence on operating cost, centered on approximately 4,980 dollars per tonne lithium carbonate. Panel b shows the top five inputs with the greatest influence on CO2 emissions, centered on approximately 5.2 tonnes CO2-equivalent per tonne lithium carbonate. Bars extend left for optimistic and right for conservative input values, showing the magnitude of each input's effect on the total.}
        \includegraphics[width=1\linewidth]{figures/Chapter 4/04-Silver Peak-Sensitivity.png}
    \altfigend
    \caption{Tornado sensitivity plots for operating costs (a) and carbon dioxide emissions (b) for the Silver Peak Project.}
    \small{The top five inputs with the greatest influence on each outcome are shown. The central vertical line indicates the midpoint estimate ($\sim$4,980/t \ce{Li2CO3} and $\sim$5.2 \ce{t CO2eq}/t \ce{Li2CO3}, respectively). Bars extend left and right in optimistic and conservative using estimates of input costs or impact factors for materials and operating designs for machine block models, respectively. Machine inputs represent design scenarios rather than continuous parameter variation, with labels representing the most significant design variation relating to the measured output; machine labels represent actual use or throughput-adjusted consumption at each scenario. Inputs whose scaled variation produced limited cost or emissions impact are not labeled.}
    \label{fig:04-silverpeak-sensitivity}
\end{figure}

The cost and emissions sensitivities shown in \ref{fig:04-silverpeak-sensitivity} tell somewhat different stories. For costs, the evaporation pond capital and operating assumptions dominate, followed by soda ash price and brine pump costs — meaning the largest uncertainty is in the pond sizing and cost model rather than reagents, even though maintenance and labor costs are relatively small (0.3-0.5\% of total capital costs as annual maintenance operations costs, 8-12 laborers managing salt removal and dike upkeep). 
For emissions, lime and soda ash embodied emissions factors dominate, followed by brine pump electricity intensity, meaning emissions uncertainty is primarily driven by upstream supply chain assumptions rather than on-site operations. This divergence is worth noting explicitly: the parameters that drive cost uncertainty are not the same as those that drive emissions uncertainty, which has implications for how policy interventions would affect the two dimensions differently.

Several modeling simplifications limit this study. Most substantively, several machine blocks (particularly thickening and dewatering, and washing and purification) are parameterized for throughputs characteristic of much larger operations (i.e. Thacker Pass and Jianxiawo); at Silver Peak's modest production scale of about 6,000 t \ce{Li2CO3} per year, these steps currently contribute negligibly to modeled costs and emissions, which may explain why our central cost estimate sits under the lower end of the literature range in \ref{fig:04-costemissions}. Additionally, other modeling simplifications are made given limited site-specific engineering data, such as incorporating pond circulation pumping costs into brine extraction pumping costs, and combining the liming pre-treatment step into the batch chemical purification step. That said, the relative consistency between our modeled range and the literature range on both cost and emissions axes, and the relative narrowness of that literature range compared to the other two projects, suggests that these simplifications do not substantially distort the overall cost and emissions competitiveness picture for a project that has been operating and reported on for nearly six decades.

With Silver Peak's relatively limited resources, its limited ability to expand due to water withdrawal limitations, and the overall lack of similar brine resources in the United States that can be exploited through the traditional evaporative pathway \cite{wesselkaemperOptimizationProcessingCosts2025}, clay resources represent a credible pathway for expanding domestic lithium production capacity, and are the next focus of this analysis. 

\subsubsection{Results for the Thacker Pass Project}\label{sec:04-thackerpass}

The Thacker Pass project, located in Humboldt County, Nevada, is a pre-production, under-construction lithium clay deposit operated by Lithium Americas Corporation. The project targets production of lithium from hectorite clay, a magnesium-rich sedimentary mineral in which lithium and magnesium occupy the same structural positions in the clay lattice \cite{castorLithiumRichClaystoneMcDermitt2020, zhaoRecentAdvancesLithium2023}. Unlike hard rock ores such as spodumene or lepidolite, which require high-temperature roasting or calcination to break down the mineral structure before lithium can be extracted \cite{gaoLithiumExtractionHard2023}, the clay structure at Thacker Pass is amenable to direct acid leaching at moderate temperatures, avoiding the most energy-intensive step of those pathways \cite{royLithiumClayAssessing2024}. However, because sulfuric acid dissolves both lithium and magnesium simultaneously, large quantities of magnesium enter the leach solution and must subsequently be removed, making magnesium management the central processing challenge of the clay pathway — analogous to its role in brine processing, but at substantially higher concentrations. Lithium Americas' latest project disclosure has defined a staged production ramp-up across four phases targeting 40,000, 80,000, 120,000, and 160,000 t \ce{Li2CO3} per year from a resource averaging 2,538 PPM lithium \cite{sgscanadaincSK1300Technical2024}. This analysis models Thacker Pass at 66,000 t \ce{Li2CO3} per year, consistent with the production volume assumed in the prior feasibility disclosure \cite{rothFeasibilityStudyNational2022}, to facilitate direct comparison with previous studies. While this is not the lifetime average across the staged production plan, we aim to represent a useful intermediate snapshot of the project's operating economics.

\begin{figure}
    \centering
    \altfig{Two-panel stacked bar chart for the Thacker Pass clay project. Panel a shows production cost in dollars per tonne lithium carbonate broken down by process step. Panel b shows greenhouse gas emissions in tonnes CO2-equivalent per tonne lithium carbonate broken down by the same steps. Error bars on each step show the sensitivity envelope from varying machine design parameters and material input costs between conservative and optimistic values.}
        \includegraphics[width=1\linewidth]{figures/Chapter 4/04-Thacker Pass-Steps-Summary.png}
    \altfigend
    \caption{Costs(a) and emissions(b) for each step of the Thacker Pass Project.}
    \small{Error bars represent a sensitivity envelope formed by simultaneously applying conservative and optimistic values for all machine design parameters and material input costs that carry defined scenario variation; they do not reflect variation in process chemistry, reagent stoichiometry, or steps for which no scenario variation is defined (e.g., tailings management). As a result, the envelope is wider for steps with more varied parameters and narrower or absent for steps with few or none. Error bars should therefore be interpreted as indicative of the dominant sources of operating uncertainty rather than as statistical confidence intervals.}
    \label{fig:04-thackerpass-summary}
\end{figure}

As described in Section~\ref{sec:04-scopemethod}, this project is modeled using a generic acid-leach clay flowsheet rather than the project-specific Thacker Pass design. The project-specific design for the Thacker Pass project incorporates a magnesium crystallization route that enables recycling of magnesium compounds as reagents, an important optimization given the magnesium-rich composition of the ore body, at the cost of additional energy usage. Excluding this route renders our estimate to be somewhat higher than the disclosure, given our central estimate of $\sim$\$8,750/t \ce{Li2CO3} compared to the full life-of-mine cost of \$8,039/t \ce{Li2CO3} reported in the disclosure \cite{sgscanadaincSK1300Technical2024}, but much higher than other estimates reported in the literature (\$6,300-\$6,986) or the 2022 disclosure (\$7,198). Notably, the first 25 years of operating costs, when ore grades are the highest, are estimated at \$6,743 and \$6,238 in the 2022 and 2024 disclosures, respectively; since our analysis did not account for ore grade degradation, our estimate is likely still an overestimate. This gap is attributable primarily to lime consumption: the generic flowsheet fully precipitates magnesium with quicklime, consuming approximately 9,200 kg \ce{CaO}/t \ce{Li2CO3}, compared to 2,630 kg/t in the project-specific design where the magnesium crystallization route recycles magnesium compounds as reagents instead. Refining the model to incorporate this route represents a priority for future work.

The step-level cost breakdown (\ref{fig:04-thackerpass-summary}a) confirms that costs are concentrated in the hydrometallurgical processing circuit rather than mining: acid leaching, impurity removal, and lithium carbonate precipitation together account for approximately \$6,100/t \ce{Li2CO3} in the midpoint baseline - over two thirds of the overall costs - with material costs dominating all three steps. Tailings disposal is the fourth most important category at $\sim$\$680/t, reflecting the large volumes of leached clay residue generated by the acid-leach process; with reagents producing these waste volumes structurally linking to handling costs in this pathway, mining operations are comparatively minor, suggesting cost reduction efforts are most productively focused on the processing circuit and its reagent use. 

The emissions breakdown (\ref{fig:04-thackerpass-summary}b) mirrors this concentration, with most emissions coming from the embodied emissions of reagents, but with a difference in ordering: impurity removal is the single largest emissions contributor at $\sim$11 t \ce{CO2eq}/t \ce{Li2CO3}, exceeding acid leaching ($\sim$2.5 t \ce{CO2eq}/t) despite ranking second in cost, reflecting the high embodied emissions of impurity-removal reagents, primarily lime. Scope 2 emissions appear in electricity-intensive mechanical steps such as crushing and attrition scrubbing, while Scope 1 emissions are largely confined to mining operations (diesel for vehicles), but these are far more minor in scope. This emissions structure implies that decarbonization of the Thacker Pass pathway depends primarily on reducing the embodied emissions of key reagents, especially lime and soda ash, rather than decarbonizing on-site electricity or fuel use.

\begin{figure}
    \centering
    \altfig{Two-panel tornado chart for the Thacker Pass project. Panel a shows the top five inputs with the greatest influence on operating cost, centered on approximately 8,750 dollars per tonne lithium carbonate. Panel b shows the top five inputs with the greatest influence on CO2 emissions, centered on approximately 17.0 tonnes CO2-equivalent per tonne lithium carbonate. Bars extend left for optimistic and right for conservative input values, showing the magnitude of each input's effect on the total.}
        \includegraphics[width=1\linewidth]{figures/Chapter 4/04-Thacker Pass-Sensitivity.png}
    \altfigend
    \caption{Tornado sensitivity plots for operating costs (a) and carbon dioxide emissions (b) for the Thacker Pass Project.}
    \small{The top five inputs with the greatest influence on each outcome are shown. The central vertical line indicates the midpoint estimate ($\sim$8,750/t \ce{Li2CO3} and $\sim$17.0 t \ce{CO2eq}/t \ce{Li2CO3}, respectively). Bars extend left and right in optimistic and conservative using estimates of input costs or impact factors for materials and operating designs for machine block models, respectively. Machine inputs represent design scenarios rather than continuous parameter variation, with labels representing the most significant design variation relating to the measured output; machine labels represent actual use or throughput-adjusted consumption at each scenario. Inputs whose scaled variation produced limited cost or emissions impact are not labeled.}
    \label{fig:04-thackerpass-sensitivity}
\end{figure}

The dominant source of cost uncertainty in the Thacker Pass Project is sulfuric acid (\ref{fig:04-thackerpass-sensitivity}a). Based on project disclosures \cite{rothFeasibilityStudyNational2022, sgscanadaincSK1300Technical2024}, we model an optimistic case of \$88/t \ce{H2SO4} reflecting on-site production costs, a midpoint of \$120/t reflecting anticipated increases in raw sulfur feedstock costs, and a conservative upper bound of \$220/t corresponding to current spot market prices, yielding a total cost range of $\sim$ \$8,100---\$10,800/t \ce{Li2CO3}. To partially mitigate this exposure, the Thacker Pass project includes an on-site sulfuric acid plant, which not only reduces both costs and price volatility relative to market procurement, but also generates additional electricity due to the exothermic nature of that production process. This plant was not explicitly modeled in this simplified generic spreadsheet. Soda ash and lime introduce secondary cost uncertainty of similar magnitude to one another, though both are smaller in range than sulfuric acid.

Of particular note is the influence of process water costs, which represent the second largest cost driver overall. Process water in Nevada is relatively expensive compared to other mining regions, and our assumed range of \$0.50---\$5.00/\si{m^3} produces meaningful cost variation at the volumetric consumption implied by the generic flowsheet. Critically, the generic flowsheet does not yet track or recycle water, which likely overstates consumption relative to the disclosed design: our model reports 16.2 million \si{m^3/yr} at 66,000 \si{t/yr} production, compared to 16.5 million \si{m^3/yr} at 160,000 \si{t/yr} in the Phase 4---5 disclosure \cite{sgscanadaincSK1300Technical2024}, a factor of approximately 2.4$\times$ higher on a per-unit basis. Given local opposition to the project's water usage \cite{alonzoenvironmentQuestionsWaterRights2025}, explicit modeling of water recycling loops represents both a technical priority and a necessary refinement before drawing conclusions about Thacker Pass's water intensity relative to the disclosed design.

From an emissions perspective, lime overwhelmingly dominates emissions uncertainties, with a range between $\sim$13.8–21.3 t \ce{CO2eq}/t \ce{Li2CO3} (\ref{fig:04-thackerpass-sensitivity}b). Because approximately two-thirds of lime's embodied emissions are process emissions chemically unavoidable from calcination \cite{lavegliaDecarbonizingLimebasedConstruction2023}, reducing the clay pathway's emissions exposure to lime requires either substituting lime in the impurity removal step with process innovation (e.g. the magnesium recycling process described above) or sourcing lime from producers deploying carbon capture, rather than simply decarbonizing grid electricity. This makes the clay pathway's emissions trajectory dependent on upstream industrial decarbonization decisions in the lime sector in a way that is structurally distinct from the electricity-dominated emissions profile of the brine pathway.

Several other modeling simplifications are also worthy of note as limitations to this study. Water recycling is not currently modeled, particularly from tailings handling and lithium carbonate production: the tailings model uses generalized dewatering cost curves rather than the project-specific lined facility described in the disclosure, which is designed to capture and recycle process water, and the lithium carbonate precipitation step likely overstates water losses relative to a design with explicit recycling loops. Furthermore, ore is reported to be slurried on-site prior to transport to the processing facility in the latest disclosure \cite{sgscanadaincSK1300Technical2024}; we instead model slurrying after the Crushing step, in the Attrition Scrubbing step. This omission likely understates both transport costs and the water balance at the mine face. 

Taken together, the generic flowsheet used here is conservative relative to the disclosed Thacker Pass design in both cost and emissions, primarily due to the absence of the magnesium crystallization route and on-site sulfuric acid production. The choice to incorporate these production steps into the Thacker Pass-specific flowsheet appears to be beneficial not only in terms of economics, but also carbon emission externalities. Despite these limitations, the modeled results suggest that Thacker Pass remains broadly competitive with Chinese lepidolite production even under conservative assumptions --- a finding that is likely to strengthen as the model is refined toward the project-specific design.

\subsubsection{Results for the Jianxiawo Project}\label{sec:04-jianxiawo}

The Jianxiawo lepidolite project, located in Yichun, Jiangxi Province, is operated by Contemporary Amperex Technology Co. Limited (CATL), currently the world's largest battery manufacturer. As China's largest lithium source, Jianxiawo provides around 3\% of global supply and 14\% of China's supply \cite{liChinasCATLSuspends2025, andrezejewskiResourcesEnergyQuarterly2025, usgeologicalsurveyMineralCommoditySummaries2026}, making it the single most consequential source of marginal lithium supply on the global cost curve.

Lepidolite is a fluorine-rich lithium mica mineral with a low bulk lithium grade, with Jiangxi lepidolite ore at approximately 0.27---0.4\% \ce{Li2O} \cite{sulcaLithiumPricesWhere2024}, or roughly 1000 to 1850 PPM Li.  Lepidolite's complex mineralogy requires high-temperature sulfate roasting to liberate the lithium. The high fluorine content in lepidolite can also generate hydrofluoric acid as a byproduct during roasting, necessitating additional reagents (e.g. lime) and generating hazardous waste streams that contribute to both costs and environmental liabilities \cite{yanExtractionLithiumLepidolite2012, dengSystematicReviewSulfate2026}. But this low grade is precisely what has allowed lepidolite projects to rapidly expand in recent years: mines operating under local kaolinite permits were not required to register as lithium mines as long as stated ore grades remained below 0.4\% \ce{Li2O}\cite{zhangChinaHasDomestic2025}. Our model deliberately assumes the upper bound of this range (0.4\%~\ce{Li2O}), to represent the most favorable ore grade at which these mines could operate under the kaolinite permit threshold, and therefore a conservative estimate of lepidolite's cost disadvantage relative to US sources.

CATL's vertical integration, which spans lithium mining through cell manufacturing, is a critical context for interpreting Jianxiawo's economics: the mine has continued operating during periods when lepidolite production costs exceeded spot lithium prices because downstream refining, cathode active material, and battery segments remained profitable \cite{jamasmieCATLCutLithium2024}. Yet the decision to idle the mine between September 2024 and February 2025, despite this integrated structure\cite{benchmarkteamCATLSuspendsOperations}, suggests that lepidolite's standalone production costs exceeded the internal transfer price at which CATL's downstream operations could profitably absorb the material — providing indirect evidence that the operation is genuinely near marginal. The standalone costs modeled here can therefore be interpreted as a meaningful benchmark for the cost at which lepidolite production enters and exits the global supply curve.

\begin{figure}
    \centering
    \altfig{Two-panel stacked bar chart for the Jianxiawo lepidolite project. Panel a shows production cost in dollars per tonne lithium carbonate broken down by process step. Panel b shows greenhouse gas emissions in tonnes CO2-equivalent per tonne lithium carbonate broken down by the same steps. Error bars on each step show the sensitivity envelope from varying machine design parameters and material input costs between conservative and optimistic values. A red vertical marker on the tailings bar indicates the modeled cost under a hazardous waste non-compliance scenario.}
        \includegraphics[width=1\linewidth]{figures/Chapter 4/04-Jianxiawo-Steps-Summary.png}
    \altfigend
    \caption{Costs(a) and emissions(b) for each step of the Jianxiawo Project.}
    \small{Error bars represent a sensitivity envelope formed by simultaneously applying conservative and optimistic values for all machine design parameters and material input costs that carry defined scenario variation; they do not reflect variation in process chemistry, reagent stoichiometry, or steps for which no scenario variation is defined (e.g., tailings management). As a result, the envelope is wider for steps with more varied parameters and narrower or absent for steps with few or none. Error bars should therefore be interpreted as indicative of the dominant sources of operating uncertainty rather than as statistical confidence intervals. The red vertical marker on the tailings bar indicates the modeled tailings cost under a hazardous waste non-compliance scenario (see Section~\ref{sec:04-jianxiawo})}. 
    \label{fig:04-jianxiawo-summary}
\end{figure}

The step-level cost breakdown (\ref{fig:04-jianxiawo-summary}a) is defined by a single dominant feature: the sulfate roasting step accounts for approximately \$5,000/t \ce{Li2CO3}, comprising material costs (primarily sodium sulfate), utility costs from high-temperature kiln operation (primarily natural gas combustion), and maintenance costs for the capital intensive, corrosion-prone roasting train, altogether representing nearly half of total modeled costs. The second largest cost is tailings disposal at $\sim$\$2,900/t, reflecting the classification of impurity sludge as hazardous waste under Chinese regulations and its associated secure landfill disposal costs \cite{nationaldevelopmentandreformcommission2024}. Outside of these two major categories, lithium carbonate precipitation contributes a further $\sim$\$900/t in material costs (soda ash), while all other steps, including mining, are comparatively minor. 

To assess the sensitivity of this cost to regulatory compliance, we modeled a scenario in which impurity sludge is handled as non-hazardous waste, i.e. non-compliance with Chinese hazardous waste classification requirements. Under this scenario, tailings disposal costs fall to approximately \$2,050/t, a reduction of roughly one-third. Thus, even under non-compliance, tailings therefore remain a substantial cost driver, because the flotation and classification tailings streams are generated at very large volumes relative to product output, in spite of low direct unit disposal costs. This tailings cost is notable in the context of the prevailing assumption that Chinese environmental regulations permit structurally lower production costs than Western jurisdictions: rather than conferring a cost advantage, hazardous waste compliance at Jianxiawo adds a cost burden largely absent from the US projects modeled here. If the cost of tailings were somehow eliminated entirely, Jianxiawo's costs could possibly undercut Thacker Pass's costs, even under the latter's project-specific magnesium recycling and on-site sulfuric acid route, but the fundamentals of processing low-grade ores and the waste materials set a floor that no regulatory environment can fully offset. 

The emissions breakdown (\ref{fig:04-jianxiawo-summary}b) is similarly dominated by the roasting step at $\sim$14.8 t \ce{CO2eq}/t \ce{Li2CO3}, but with an important structural distinction from the Thacker Pass pathway: roasting emissions are divided roughly equally between Scope 1 direct combustion emissions from high-temperature kiln operation and Scope 3 embodied emissions from sodium sulfate and lime reagents \cite{yanExtractionLithiumLepidolite2012, tranChapter3Lithium2015}. This means that even full decarbonization of the natural gas combustion component would leave roughly half of roasting emissions unaddressed, and that emissions reductions require both fuel switching and upstream reagent supply chain decarbonization simultaneously. Outside of the crucial sulfate roasting step, lithium carbonate precipitation is the only other notable emissions contributor at $\sim$4.8 t \ce{CO2eq}/t, driven by the embodied emissions of soda ash. It is important to note that this step likely has a higher environmental impact in China than in the United States, given that the soda ash used in this step is likely produced from the heat-intensive Solvay process ($\sim$0.7-1.2 t \ce{CO2eq}/t \ce{Na2CO3}) \cite{wangGreenhouseGasesRegulated2024, scottCanSyntheticSoda2023, royLithiumClayAssessing2024}. All other steps have much smaller impacts in comparison.

\begin{figure}
    \centering
    \altfig{Two-panel tornado chart for the Jianxiawo project. Panel a shows the top five inputs with the greatest influence on operating cost, centered on approximately 10,400 dollars per tonne lithium carbonate. Panel b shows the top five inputs with the greatest influence on CO2 emissions, centered on approximately 22.6 tonnes CO2-equivalent per tonne lithium carbonate. Bars extend left for optimistic and right for conservative input values, showing the magnitude of each input's effect on the total.}
        \includegraphics[width=1\linewidth]{figures/Chapter 4/04-Jianxiawo-Sensitivity.png}
    \altfigend
    \caption{Tornado sensitivity plots for operating costs (a) and carbon dioxide emissions (b) for the Jianxiawo Project.}
    \small{The top five inputs with the greatest influence on each outcome are shown. The central vertical line indicates the midpoint estimate ($\sim$10,400/t \ce{Li2CO3} and $\sim$22.6 \ce{t CO2eq}/t \ce{Li2CO3}, respectively). Bars extend left and right in optimistic and conservative using estimates of input costs or impact factors for materials and operating designs for machine block models, respectively. Machine inputs represent design scenarios rather than continuous parameter variation, with labels representing the most significant design variation relating to the measured output; machine labels represent actual use or throughput-adjusted consumption at each scenario. Inputs whose scaled variation produced limited cost or emissions impact are not labeled.}
    \label{fig:04-jianxiawo-sensitivity}
\end{figure}

The cost sensitivity for the Jianxiawo lepidolite pathway (\ref{fig:04-jianxiawo-sensitivity}a) is shaped by three large contributing variables. Soda ash price produces the widest range at approximately \$1,100/t, reflecting both the large stoichiometric requirement of soda ash for lithium carbonate precipitation and the wide spread between Chinese Solvay-process and US trona-derived soda ash prices. Sodium sulfate price and the sulfate roasting train design (the latter driven primarily by varying combustion efficiency assumptions in the roasting process) follow at approximately \$900/t and \$800/t, respectively. This is structurally different from Thacker Pass, where a single reagent dominated, and implies that cost reduction at Jianxiawo requires simultaneous optimization across the roasting circuit and multiple reagent supply chains rather than a single-lever intervention. Lime and process water contribute secondary uncertainty of smaller magnitudes. 

In contrast to the cost sensitivity, where uncertainty is distributed roughly equally across three drivers, emissions uncertainty is dominated by a single factor: sodium sulfate embodied emissions span $\sim$20---28 t \ce{CO2eq}/t \ce{Li2CO3} (\ref{fig:04-jianxiawo-sensitivity}b). This range is notably wider than the next three largest drivers: the roasting machine energy intensity, soda ash, and lime follow with roughly equal ranges of $\sim$4 t \ce{CO2eq}/t, consistent with the step-level findings discussed above.

Several limitations specific to the Jianxiawo model warrant elaboration. Most fundamentally, lepidolite processing is considerably less documented in the open literature than either brine or clay pathways: the sulfate roasting flowsheet is reconstructed from patents filed in China and lab-scale studies \cite{yanExtractionLithiumLepidolite2012, luongFactorsAffectingExtraction2013, parkEffectiveLithiumExtraction2023, dengSystematicReviewSulfate2026, nanMethodsExtractingLithium2022, nanMethodsExtractingLithium2019, zouMethodEfficientLithium2024} rather than from engineering disclosures or feasibility studies comparable to those available for Thacker Pass and Silver Peak, which are not available given the lack of disclosure requirements for Chinese-listed mining companies. This data gap is itself a small finding, reflecting the limited public transparency of Chinese lithium production operations, but also means that process parameters such as reagent ratios, roasting temperatures, and conversion efficiencies carry greater uncertainty than for the US projects, and that the modeled flowsheet may not accurately represent CATL's proprietary process design. Similarly, the tailings model is relatively simple: impurity sludge is assigned a hazardous disposal cost based on published Chinese regulatory gate fees \cite{nationaldevelopmentandreformcommission2024}, but the volume and composition of sludge generated are derived from the generic flowsheet rather than measured plant data, introducing uncertainty in what is the second largest cost category. Additionally, the production volume modeled here, based on the most commonly reported production figure of 46,000 tonnes LCE, is less than half of the highest reported capacity of 111,000 t LCE \cite{benchmarkteamRestartCATLsJianxiawo2026, andrezejewskiResourcesEnergyQuarterly2025}, and thus carries more uncertainty than projects owned by US or Canadian-listed companies, which are anchored to stock market-required engineering disclosures. Given the presence of economies of scale in the roasting circuit, cost and emissions estimates are sensitive to this assumption in ways that are difficult to bound without plant-level data. Finally, as noted above, the standalone costs modeled here do not capture the cost-sharing or transfer-pricing arrangements that may apply within CATL's vertically integrated operations, and thus may not fully reflect the true `marginal' lithium producer in China.

Our bottom-up model suggests that neither the cost nor the emissions profile of Jianxiawo is competitive with Silver Peak, and the overlap with Thacker Pass in the cost-emissions space is partially attributable to conservative generic flowsheet assumptions applied to the latter and the most favorable ore grade permissible under the kaolinite permit threshold (0.4\% \ce{Li2O}) for the former, suggesting the true gap between US clay and Chinese lepidolite production is likely wider than our midpoint estimates indicate. At the same time, our process-based estimates fall below figures reported in the literature and industry reports \cite{gaoLithiumExtractionHard2023, atkinsonLithiumPriceReversion2023, liuChinaLithiumBoom2024}, suggesting that prior high estimates likely reflect data scarcity and contextual limitations. As a result, lepidolite's cost and emissions disadvantage, while real and structural, may be somewhat less severe than what has been commonly reported. Both conclusions point in the same direction: more detailed, process-grounded modeling of Chinese lepidolite operations is needed to bound these estimates, and represents a priority for future work.

\subsection{Limitations}\label{sec:04-limitations}

Several important scope limitations of the present analysis should be addressed in future work, beyond those that were discussed for each individual project. First, mine closure and reclamation costs are not modeled here. These costs can be substantial, particularly for open-pit mining operations like Thacker Pass and Jianxiawo, and their exclusion likely understates the full lifecycle cost of all three projects — though the relative magnitude of this omission may differ across jurisdictions given differences in reclamation requirements between the United States and China.

Second, this analysis reports production costs at the mine gate and does not model cost, insurance, and freight (CIF) delivery costs to end markets. Chinese producers may retain a competitive advantage in delivered cost to Asian battery manufacturers even where their production costs are higher, given proximity to demand centers and established logistics infrastructure. US producers serving domestic battery manufacturers could benefit from localized supply chains, but establishing these logistics channels from rural Nevada to battery manufacturing plants primarily located in the US automotive corridor \cite{chengUSIndustrialPolicy2024} may be more costly in comparison to China. A full delivered-cost analysis, incorporating transportation, tariffs, and insurance, would provide a more complete picture of competitive positioning under different trade policy scenarios.

Further, water consumption is an important dimension that has only been partially addressed in this analysis. While a water recycling loop was incorporated into the model for Silver Peak, where solar evaporation dominates water consumption, water flows are not yet explicitly tracked or recycled in the Thacker Pass or Jianxiawo models. Explicit water balance modeling for all three projects is a priority refinement, both for accuracy and because water usage is increasingly a determinant of the local social license to operate, especially in water-stressed regions such as Nevada \cite{sunRobustAssessmentsLithium2025, schombergExtendedLifeCycle2021}.

Additionally, tailings and residue management costs are modeled using a simplified piece-wise empirical dewatering cost function, parameterized by solids fraction, combined with project-specific disposal routes~\cite{sgscanadaincSK1300Technical2024,srkconsultingSECTechnicalReport2023a} and gate fees for hazardous and non-hazardous waste~\cite{nationaldevelopmentandreformcommission2024}. While the dominant cost drivers are directly observed, particularly the hazardous disposal burden at Jianxiawo, which represents the paper's largest single finding on environmental compliance costs, some parameters within the model remain as order-of-magnitude placeholders, including on-site tailings facility operating costs and wastewater treatment costs. These are unlikely to affect the paper's comparative conclusions materially, given that disposal costs at the US projects are small relative to reagent costs, but represent a priority refinement for site-specific validation. More substantively, the environmental impacts of waste disposal are not currently modeled. For the ore-based pathways, where large volumes of leach residue and impurity sludge are generated and need to be dewatered, this omission likely understates total Scope 1 and 2 emissions modestly, though the dominant emissions drivers identified in the results (reagent embodied emissions and fuel combustion) are unaffected by this gap.

Beyond these refinements to existing models, several extensions would broaden the contribution of this work. The three projects modeled here are representative but not comprehensive: several other clay projects exist in the United States at various stages of development \citep{wesselkaemperOptimizationProcessingCosts2025}, and other lepidolite and spodumene operations in China could be modeled. In particular, the brine-direct lithium extraction (DLE) pathway, with projects under development in both the United States and China, appears unlikely to be marginal at current prices, but would be important for characterizing the full domestic cost curves of both countries, and understanding how the competitive landscape shifts under higher-price scenarios. Altogether, these additions would enable a richer analysis of how cost curves shift under different demand growth scenarios and how individual projects enter and exit production as prices move, providing an empirical foundation for a more complete treatment of supply curve dynamics.

All of these extensions described above, like this current analysis, are static in nature, characterizing production costs and emissions under a given set of assumptions and uncertainties. However, with project-level cost curves established for a broader set of pathways, they can serve as inputs to dynamic models that can capture firm- and market-level responses to policy interventions, technological learning, and shifting market conditions. The progression from static project baselines to dynamic supply curve modeling is therefore a natural next step, and the framework developed here is designed with that extension in mind.

\subsection{Conclusions, Policy Implications, Future Work}

This study developed bottom-up technoeconomic-life cycle assessment models for three lithium production projects representing distinct resource-pathway combinations: Silver Peak (evaporative brine, United States), Thacker Pass (sulfuric acid-leach clay, United States), and Jianxiawo (sulfate roasting lepidolite, China). The results demonstrate that the cost and emissions profiles of lithium production are not separable from the resource from which it is extracted: the thermochemical and hydrometallurgical requirements of each pathway create structural differences in production costs, emissions intensities, and sensitivity to price volatility that are challenging to overcome with project-specific optimizations. Under central assumptions, both US projects produce lithium carbonate at a lower cost and life cycle emissions than the Chinese lepidolite operation that appears to set the marginal price on the global supply curve, with Silver Peak being distinctly favorable in both dimensions, and Thacker Pass remaining broadly competitive even under conservative generic flowsheet assumptions. Crucially, while Silver Peak has limited ability to expand due to limited land and water rights, the opening of the Thacker Pass project, which when operating at maximum production capacity could supply 10\% of current global lithium demand \cite{usgeologicalsurveyMineralCommoditySummaries2026}, could establish the United States as a major producer of lithium as demand continues to grow. 

The United States holds some structural natural resource advantages relevant to lithium production costs and emissions, though these should not be overstated given China's extensive and low-cost chemical processing industry. One example is soda ash: US producers source it from mined trona in Wyoming's Green River Basin, the world's largest natural soda ash deposit \cite{usgeologicalsurveyMineralCommoditySummaries2026}, at substantially lower cost and emissions intensity than the Solvay process predominant in China \cite{wangGreenhouseGasesRegulated2024, scottCanSyntheticSoda2023, royLithiumClayAssessing2024}, conferring a structural advantage across all three pathways modeled here. More broadly, the abundant solar and wind resources of the arid Southwest where US lithium deposits are concentrated, and relatively low natural gas prices represent potential further advantages as energy costs become more significant at larger production scales. 

On the other hand, a common assumption in supply chain security discourse is that Chinese production benefits from lower environmental regulatory costs. Our results complicate this narrative: hazardous waste disposal requirements for lepidolite roasting residues contribute approximately \$2,900/t \ce{Li2CO3} (more than 25\%) to Jianxiawo's modeled costs — a burden largely absent from the US projects, which remains notable even under a hypothetical non-compliance scenario in which hazardous tailings were treated as non-hazardous (\$2,050/t \ce{Li2CO3}). Furthermore, there is evidence that Chinese environmental compliance costs for lepidolite operations are likely to increase rather than decrease: as noted in Section~\ref{sec:04-jianxiawo}, many lepidolite mining permits in Jiangxi Province were originally issued as kaolinite permits, and ongoing regulatory tightening around this permitting gap \cite{benchmarkteamCATLSuspendsOperations, zhangChinaHasDomestic2025} could raise operating costs and restrict capacity in ways not captured in current cost curve estimates.

These project-level baselines also could provide the foundation for a systematic analysis of policy interventions. Building on the models developed here, future work could analyze the relative competitiveness of each pathway under existing policies such as including production tax credits under the IRA and OBBBA, Made in China 2025 industrial subsidies, and Chinese export restrictions, or prospective interventions such as a carbon border adjustment mechanism (CBAM), currently being implemented in Europe. Such analyses would allow policymakers to identify which interventions most efficiently close the competitiveness gap between US and Chinese production, and under what conditions domestic lithium supply development becomes self-sustaining without ongoing policy support.

More broadly, this framework for integrated process-level technoeconomic and life cycle emissions assessment across geographic production pathways does not need to be specific to lithium. While the specific projects and economic landscape analyzed here reflect a particular moment in the geopolitics of critical minerals, the framework itself is material-agnostic: wherever production costs and environmental outcomes are jointly shaped by resource type and production location, the same approach could apply. Applied to other critical mineral systems, this approach can contribute to a more rigorous and consistent basis for evaluating supply chain security, environmental performance, and the effectiveness of industrial policy across the materials and manufacturing systems underpinning the energy transition.